\begin{multicols}{2}

\chapter{Space Effects}\index{Science Fiction Rules!Space Effects}

\section{Effects of Gravity}\index{Science Fiction Rules!Space Effects!Gravity}

In a hard science fiction game, variable effects of gravity can’t be waved away by tech that simulates normal gravity on spacecraft, space stations, and other worlds. Instead, it’s an issue people must overcome. 

\textbf{Short-Term Microgravity Exposure}: People new to low gravity might get space sickness. Newcomers must succeed on a difficulty 3 Might task or suffer mild nausea for about two to four days, during which time all their tasks are hindered. A few unlucky travelers (usually those who roll a 1 or otherwise face a GM intrusion) are almost completely incapacitated, and find all tasks hindered by three steps.

\textbf{Long-Term Microgravity Exposure}: Long-term exposure to microgravity environments without medical interventions degrades health. How long one spends in such conditions is directly relevant. The GM may assign long-term penalties to PCs if the situation warrants it, though the use of advanced space medicine, proper exercise, and recommended steroids and other hormones can avoid these complications.

\textbf{Low Gravity}: Weapons that rely on weight, such as all heavy weapons, inflict 2 fewer points of damage (dealing a minimum of 1 point). Short-range weapons can reach to long range, and long-range weapons can reach to very long range. Characters trained in low-gravity maneuvering ignore the damage penalty.

\textbf{High Gravity}: It’s hard to make effective attacks when the pull of gravity is very strong. Attacks (and all physical actions) made in high gravity are hindered. Ranges in high gravity are reduced by one category (very-long-range weapons reach only to long range, long-range weapons reach only to short range, and
short-range weapons reach only to immediate range). Characters trained in high-gravity maneuvering ignore the change in difficulty but not the range decreases.

\textbf{Zero Gravity}: It’s hard to maneuver in an environment without gravity. Attacks (and all physical actions) made in zero gravity are hindered. Short-range weapons can reach to long range, and long-range weapons can reach to very-long range.

\section{Effects of Vacuum}\index{Science Fiction Rules!Space Effects!Vacuum}

Vacuum is lethal. There’s no air to breathe, and the lack of pressure causes havoc on an organic body. An unprotected character moves one step down the damage track each round. However, at the point where they should die, they instead fall unconscious and remain so for about a minute. If they are rescued during that time, they can be revived. If not, they die.

\chapter{Travelling the Solar System and Orbital Mechanics}\index{Science Fiction!Travelling the Solar System and Orbital Mechanics}

In a hard science fiction setting, you might be interested in evoking the reality of travel times between colonies on planets and moons in the solar system. Even so, plotting a course between locations in the solar system isn’t simple, because everything is always moving with respect to everything else. You could determine exactly how long a trip would take with some internet research. Or you could just evoke the effect of orbital mechanics and varying accelerations on interplanetary travel. Use the Interplanetary Travel Table to do so. For a trip between locations not directly compared, add up the destinations in between. The travel times assume a nuclear plasma engine of a kind already being tested today (but better), a steady thrust toward the destination, and an equally long and steady braking thrust over the last half of the trip before orbit insertion. Such propulsion systems can change velocity and sustain thrust for days at a time, which reduces bone loss, muscle atrophy, and other long-term effects of low gravity.

Regardless, the travel times between distant locations bring home one thing: space is big and lonely.

\end{multicols}

\begin{center} 
\includegraphics{d20div}
\end{center}
